\section*{الفصل \ref{ch:b2:catalytic-cycles} --- الحفز العضوي الفلزي: الخطوات الأولية والدورات}

\begin{solution}{exo:b2:catalytic-cycles:1}
قبل: Ir(I)، $d^8$، $8 + 4 \times 2 = 16$ إلكترونًا. بعد: Ir(III)، $d^6$، $6 + 6 \times 2 = 18$
إلكترونًا، ثماني الأوجه.
\end{solution}

\begin{solution}{exo:b2:catalytic-cycles:2}
الأولى \omterm{def:b2:catalytic-cycles:oxidative-addition}{حذف اختزالي} (من Pt(IV) إلى Pt(II)، ويتكوّن الإيثان من مجموعتي ميثيل).
والثانية \omterm{def:b2:catalytic-cycles:ligand-exchange}{تبادل ربائط} (يحل الإيثين محل فوسفين واحد) يتبعه
\omterm{def:b2:catalytic-cycles:insertion}{إدخال هجري} للإيثين في الرابطة Pd--Ph.
\end{solution}

\begin{solution}{exo:b2:catalytic-cycles:3}
\babelsublr{TON $= 15.0/0.020 = 750$}؛ \babelsublr{TOF $= 750/3.0 = \qty{250}{h^{-1}}$}.
\end{solution}

\begin{solution}{exo:b2:catalytic-cycles:4}
\begin{align*}
  &\mathrm{PdL_2 + ArBr \to ArPdBrL_2}\\
  &\mathrm{ArPdBrL_2 + Ar'B(OH)_3^- \to ArPdAr'L_2 + Br^- + B(OH)_3}\\
  &\mathrm{ArPdAr'L_2 \to Ar{-}Ar' + PdL_2}
\end{align*}
المجموع: $\mathrm{ArBr + Ar'B(OH)_3^- \to Ar{-}Ar' + Br^- + B(OH)_3}$؛ وكل أنواع البلاديوم
تُختزل.
\end{solution}

\begin{solution}{exo:b2:catalytic-cycles:5}
الإيثيل والإيزوبروبيل (لهما هيدروجينات $\beta$). فالميثيل ليس له كربون $\beta$؛ والكربون $\beta$
في النيوبنتيل لا يحمل هيدروجينًا؛ والكربون $\beta$ في البنزيل كربون حلقي \omterm{def:b2:huckel:aromatic}{عطري} لا يحمل
هيدروجينًا.
\end{solution}

\begin{solution}{exo:b2:catalytic-cycles:6}
الناتج (\textit{E})-3-فينيل بروبينوات الميثيل، \ce{Ph-CH=CH-COOCH3}، والفينيل مضاف إلى
الكربون الطرفي، والناتج trans.
\end{solution}

\begin{solution}{exo:b2:catalytic-cycles:7}
البيوتانال \ce{CH3CH2CH2CHO} (خطي: يُضاف الفلز إلى الكربون الطرفي، والهيدريد إلى
الكربون الداخلي) و2-ميثيل بروبانال \ce{(CH3)2CHCHO} (متفرع: الفلز على الكربون
الداخلي).
\end{solution}

\begin{solution}{exo:b2:catalytic-cycles:8}
للمعقد \ce{[RhCl(PPh3)3]} ستة عشر إلكترونًا: فإضافة \ce{H2} (+2) تعطي 18 مع ستة
مواضع تناسق، وهو ازدحام شديد مع ثلاثة فوسفينات ضخمة. أما \ce{[RhCl(PPh3)2]} فله أربعة عشر
وموقع مفتوح: فتعطي \omterm{def:b2:catalytic-cycles:oxidative-addition}{الإضافة التأكسدية} بسهولة معقدًا ذا 16 إلكترونًا.
\end{solution}

\begin{solution}{exo:b2:catalytic-cycles:9}
حمض البورونيك محب للنواة ضعيف؛ وتضيف القاعدة الهيدروكسيد (أو الألكوكسيد) إلى البور،
فتعطي بورات \ce{ArB(OH)3-}، مجموعتها الأريلية أغنى بالإلكترونات وتنتقل إلى
البلاديوم.
\end{solution}

\begin{solution}{exo:b2:catalytic-cycles:10}
المعقد ذو العشرين إلكترونًا غير معقول. فيجب أن يفقد المعقد أولًا ربيطة (14 إلكترونًا)، ثم
يضيف \ce{H2} (16)، ويرتبط بالألكين (18)، ويدخله في رابطة M--H (16) ويحذف
الألكان (14)، قبل أن يستعيد الربيطة أو يبدأ من جديد.
\end{solution}

\begin{solution}{exo:b2:catalytic-cycles:11}
السرعة الابتدائية $n_{\max}/\tau = \qty{0.25}{mmol/min}$؛ وTOF الابتدائي $= 0.25/0.010 =
\qty{25}{min^{-1}} = \qty{1.5e3}{h^{-1}}$. وTON النهائي $= 10.0/0.010 = 1.0 \times 10^3$.
\end{solution}

\begin{solution}{exo:b2:catalytic-cycles:12}
4-بروموأنيسول مع حمض الفينيل بورونيك (أو بروموبنزين مع حمض 4-ميثوكسي فينيل بورونيك)،
وحفاز فوسفيني للبلاديوم(0) وقاعدة. الدورة: \omterm{def:b2:catalytic-cycles:oxidative-addition}{الإضافة التأكسدية}
لبروميد الأريل، ثم \omterm{def:b2:catalytic-cycles:transmetalation}{التبادل الفلزي} من البورات، ثم \omterm{def:b2:catalytic-cycles:oxidative-addition}{الحذف الاختزالي} للمركب
4-ميثوكسي ثنائي الفينيل.
\end{solution}

\begin{solution}{pb:b2:catalytic-cycles:1}
\textbf{1.} Pd(II)، $d^8$.
\textbf{2.} Pd(0)، $d^{10}$، $10 + 2 \times 2 = 14$ إلكترونًا.
\textbf{3.} بأربعة عشر إلكترونًا يمكنه استقبال زوجين آخرين: فله مواقع شاغرة.
\textbf{4.} معقدات البلاديوم(0) الفوسفينية يؤكسدها الهواء، والفوسفينات
تتأكسد إلى أكاسيد الفوسفين.
\textbf{5.} يختزل البلاديوم(II) إلى البلاديوم(0) (متأكسدًا هو نفسه) ويرتبط
بالبلاديوم(0)، فيبقيه في المحلول.
\textbf{6.} \omterm{def:b2:catalytic-cycles:cycle}{طليعة الحفاز}.
\textbf{7.} \ce{Pd(PPh3)2 + CH3C6H4Br -> [Pd(C6H4CH3)Br(PPh3)2]}: Pd(II)، $d^8$، $8 + 4 \times 2 =
16$ إلكترونًا.
\textbf{8.} تحوّله إلى البورات \ce{PhB(OH)3-}، التي تنقل مجموعتها الفينيلية.
\textbf{9.} \ce{[Pd(Ar)Br(PPh3)2] + PhB(OH)3- -> [Pd(Ar)(Ph)(PPh3)2] + Br- + B(OH)3}.
\textbf{10.} \omterm{def:b2:catalytic-cycles:oxidative-addition}{الحذف الاختزالي} يصل ربيطتين في وضع cis؛ فمجموعتا الأريل في وضع trans يجب أن تتماكبا أولًا.
\textbf{11.} \ce{[Pd(Ar)(Ph)(PPh3)2] -> Ar-Ph + Pd(PPh3)2}: بلاديوم(0)، 14 إلكترونًا، جاهز
لدورة أخرى.
\textbf{12.} \ce{CH3C6H4Br + PhB(OH)3- -> CH3C6H4-Ph + Br- + B(OH)3}.
\textbf{13.} 4-بروموتولوين (\qty{10.0}{mmol}، مقابل \qty{12.0}{mmol} من حمض البورونيك).
\textbf{14.} $1.51/168.24 = \qty{8.98e-3}{mol}$، أي \qty{8.98}{mmol}.
\textbf{15.} $8.98/10.0 = 89.8~\%$.
\textbf{16.} $0.050/10.0 = 0.50$~mol~\%.
\textbf{17.} $8.98/0.050 = 180$.
\textbf{18.} $180/4.0 = \qty{45}{h^{-1}}$.
\textbf{19.} مجموعتا فينيل انتقلتا إلى ذرة البلاديوم نفسها (الاقتران المتجانس لحمض
البورونيك، الذي تيسّره آثار الأكسجين)، ثم حُذفتا معًا.
\textbf{20.} ليس لمجموعات الأريل هيدروجين $\beta$ قادر على بلوغ الفلز.
\textbf{21.} \omterm{def:b2:catalytic-cycles:oxidative-addition}{الإضافة التأكسدية}: فالرابطة \ce{C-Cl} أقوى من \ce{C-Br}.
\textbf{22.} ينتهي جسيماتٍ من البلاديوم(0) أو أملاحًا في البقايا؛ وهو باهظ الثمن 
وسام، فيُستعاد.
\textbf{23.} يضيع جزء من حمض البورونيك في تفاعلات جانبية (الاقتران المتجانس، وفقد البور
من الحلقة)؛ والفائض يُبقي البروميد هو المحدد.
\textbf{24.} دار البلاديوم $\boldsymbol{\approx \num{1.8e2}}$ مرة.
\end{solution}
